\section{Resumen}

Esta tesis estudia el seguimiento de precios en horizonte finito en un creador automático de mercado (AMM) sintético, de producto constante y sin comisiones, frente a un choque en una referencia exógena de un mercado centralizado (CEX). Separa la mecánica exacta del intercambio de los efectos de la capacidad de arbitraje, la profundidad de reservas, las reglas de los agentes y el orden de eventos. En una piscina aislada, la capacidad $c$ corrige el error logarítmico fuera de banda hasta en $2\log(1+c)$ por tick; un choque de 0.20 con banda 0.003 requiere cuatro ticks ideales con $c=0.03$. El estudio reproducible evalúa 13,500 trayectorias principales, una grilla de sensibilidad de 40,500 y una repetición de 9,000 trayectorias archivadas.

El arbitraje de capacidad alta reduce el área media de exceso de error acumulado de 2.36 sin arbitraje a 0.10, aunque la dirección del contraste está incorporada en la regla. El retiro único se activa en todas las trayectorias E3. En la repetición poco profunda, las salidas posteriores de E3 aparecen en el 98.2--99.6\,\% de las trayectorias recuperadas y a menudo se repiten; un control con capacidad fija iguala las razones configuradas con E1 estático y reduce el contraste del área posrecuperación de 0.122--0.124 a 0.0036--0.0097. Una comprobación dirigida de 900 trayectorias con capacidad fija y sin pánico produce cero salidas posteriores en cada celda de choque, lo que respalda $R$ como hipótesis organizadora candidata, no como ley validada. El control es posterior al análisis inicial, cubre una profundidad y condiciona el resultado a la recuperación; no identifica efectos causales separados de capacidad y profundidad.

El mapa reducido de dos estados mejora frente a la persistencia en algunas condiciones reservadas, pero no frente al AR(1) escalar. Un mapa híbrido sensible a la banda reduce cerca del 1.7\,\% el error de reproducción recursiva usando la trayectoria futura realizada del CEX; por ello no es un pronóstico en tiempo real. Los resultados son condicionales al conjunto de reglas sintéticas y no tienen validación con datos de mercado. El aporte es un estudio reproducible de mecanismos que identifica telemetría y protocolos de reproducción para trabajos futuros con precios sincronizados, comisiones, demoras, costos de operación, capital finito y competencia entre arbitrajistas realistas.

\textbf{Palabras clave:} creador automático de mercado; arbitraje; seguimiento de precios; retiro de liquidez; simulación basada en agentes; Monte Carlo
